
\documentclass{article}
\usepackage{colortbl}
\usepackage{makecell}
\usepackage{multirow}
\usepackage{supertabular}

\begin{document}

\newcounter{utterance}

\centering \large Interaction Transcript for game `hot\_air\_balloon', experiment `air\_balloon\_survival\_en\_negotiation\_hard', episode 3 with glm{-}4.6{-}t1.0.
\vspace{24pt}

{ \footnotesize  \setcounter{utterance}{1}
\setlength{\tabcolsep}{0pt}
\begin{supertabular}{c@{$\;$}|p{.15\linewidth}@{}p{.15\linewidth}p{.15\linewidth}p{.15\linewidth}p{.15\linewidth}p{.15\linewidth}}
    \# & $\;$A & \multicolumn{4}{c}{Game Master} & $\;\:$B\\
    \hline

    \theutterance \stepcounter{utterance}  
    & & \multicolumn{4}{p{0.6\linewidth}}{
        \cellcolor[rgb]{0.9,0.9,0.9}{
            \makecell[{{p{\linewidth}}}]{
                \texttt{\tiny{[P1$\langle$GM]}}
                \texttt{You are participating in a collaborative negotiation game.} \\
\\ 
\texttt{Together with another participant, you must agree on a single set of items that will be kept. Each of you has your own view of how much each item matters to you (importance). You do not know how the other participant values the items. Additionally, you are given the effort each item demands.} \\
\texttt{You may only agree on a set if the total effort of the selected items does not exceed a shared limit:} \\
\\ 
\texttt{LIMIT = 3681} \\
\\ 
\texttt{Here are the individual item effort values:} \\
\\ 
\texttt{Item effort = \{"C17": 700, "B41": 586, "A71": 505, "C00": 116, "B35": 664, "A77": 940, "A03": 387, "A04": 392, "C54": 209, "C01": 571, "A87": 4, "B32": 285, "B86": 651, "B33": 613, "A10": 740\}} \\
\\ 
\texttt{Here is your personal view on the importance of each item:} \\
\\ 
\texttt{Item importance values = \{"C17": 138, "B41": 583, "A71": 868, "C00": 822, "B35": 783, "A77": 65, "A03": 262, "A04": 121, "C54": 508, "C01": 780, "A87": 461, "B32": 484, "B86": 668, "B33": 389, "A10": 808\}} \\
\\ 
\texttt{Goal:} \\
\\ 
\texttt{Your goal is to negotiate a shared set of items that benefits you as much as possible (i.e., maximizes total importance to YOU), while staying within the LIMIT. You are not required to make a PROPOSAL in every message {-} you can simply negotiate as well. All tactics are allowed!} \\
\\ 
\texttt{Interaction Protocol:} \\
\\ 
\texttt{You may only use the following structured formats in a message:} \\
\\ 
\texttt{PROPOSAL: \{'A', 'B', 'C', …\}} \\
\texttt{Propose keeping exactly those items.} \\
\texttt{REFUSE: \{'A', 'B', 'C', …\}} \\
\texttt{Explicitly reject opponent's proposal.} \\
\texttt{ARGUMENT: \{'...'\}} \\
\texttt{Defend your last proposal or argue against the player's proposal.} \\
\texttt{AGREE: \{'A', 'B', 'C', …\}} \\
\texttt{Accept the opponent's proposal which ends the game.} \\
\texttt{STRATEGIC REASONING: \{'...'\}} \\
\texttt{	Describe your strategic reasoning or anticipation explaining your choice of action. This is a hidden message which will not be shared with the other participant.} \\
\\ 
\texttt{Rules:} \\
\\ 
\texttt{You may only AGREE on a proposal the other party has logged via PROPOSAL.} \\
\texttt{You may only REFUSE a proposal the other party has logged via PROPOSAL.} \\
\texttt{Total effort of any PROPOSAL or AGREE set must be ≤ LIMIT.} \\
\texttt{Do NOT reveal your hidden importance scores.} \\
\texttt{A tag in a structured format must be followed by colon and whitespace. The argument must be a python set containing 0 or more strings.} \\
\texttt{So, it must be of the form TAG: \{...\}} \\
\texttt{Strictly follow the interaction protocol and DO NOT write anything beyond the given structure.} \\
\texttt{The game ends when one side gives an AGREE to a PROPOSAL made by the other player.} \\
\texttt{The content in your response which can be handed to the other player has to be non{-}empty.} \\
\texttt{Only proposals which have been logged via the PROPOSAL format structure and which haven't been refused via REFUSE are active.} \\
\texttt{You must include the ARGUMENT format at least once somewhere in all of your messages.} \\
\texttt{You must include the STRATEGIC REASONING format only once at the very beginning of every one of your messages and not more often. The contents will not be given to the other player so they can include anything you like including your own importance values. Here you should reason step by step to come up with you next move.} \\
            }
        }
    }
    & \\ \\

    \theutterance \stepcounter{utterance}  
    & \multicolumn{4}{p{0.6\linewidth}}{
        \cellcolor[rgb]{0.9,0.9,0.9}{
            \makecell[{{p{\linewidth}}}]{
                \texttt{\tiny{[P1$\rangle$GM]}}
                \texttt{STRATEGIC REASONING: \{'My first step is to determine my ideal set of items, the one that maximizes my personal importance score within the effort limit. This is a 0/1 knapsack problem. I calculate the importance{-}to{-}effort ratio for each item. My highest{-}value items are A87 (ratio: 115.25), C00 (ratio: 7.09), C54 (ratio: 2.43), A71 (ratio: 1.72), B32 (ratio: 1.70), C01 (ratio: 1.37), B35 (ratio: 1.18), B86 (ratio: 1.03), and B41 (ratio: 0.99). Selecting items based on this ratio gives me the set \{A87, C00, C54, A71, B32, C01, B35, B86, B41\}. The total effort for this set is 3591, which is under the limit of 3681. The total importance I derive from this set is 5957. This is my optimal outcome. My strategy will be to propose this set immediately. It offers me maximum value and leaves very little remaining effort (only 90), which restricts the opponent\textbackslash{}'s ability to add their own items and forces a negotiation based on swaps rather than additions, giving me a strong negotiating position. If they refuse, I will be prepared to trade away some of my lower{-}importance items from this set (like A87, C54, B32, or B41) for items they value that I value less.'\}} \\
\texttt{PROPOSAL: \{'A87', 'B41', 'A71', 'C00', 'B35', 'C54', 'C01', 'B32', 'B86'\}} \\
\texttt{ARGUMENT: \{'This proposal makes excellent use of our shared effort limit, selecting a strong combination of items that collectively reach an effort of 3591. It represents a highly efficient allocation of resources, ensuring we keep a diverse and valuable set of items.'\}} \\
            }
        }
    }
    & & \\ \\

    \theutterance \stepcounter{utterance}  
    & & \multicolumn{4}{p{0.6\linewidth}}{
        \cellcolor[rgb]{0.9,0.9,0.9}{
            \makecell[{{p{\linewidth}}}]{
                \texttt{\tiny{[GM$\rangle$P2]}}
                \texttt{You are participating in a collaborative negotiation game.} \\
\\ 
\texttt{Together with another participant, you must agree on a single set of items that will be kept. Each of you has your own view of how much each item matters to you (importance). You do not know how the other participant values the items. Additionally, you are given the effort each item demands.} \\
\texttt{You may only agree on a set if the total effort of the selected items does not exceed a shared limit:} \\
\\ 
\texttt{LIMIT = 3681} \\
\\ 
\texttt{Here are the individual item effort values:} \\
\\ 
\texttt{Item effort = \{"C17": 700, "B41": 586, "A71": 505, "C00": 116, "B35": 664, "A77": 940, "A03": 387, "A04": 392, "C54": 209, "C01": 571, "A87": 4, "B32": 285, "B86": 651, "B33": 613, "A10": 740\}} \\
\\ 
\texttt{Here is your personal view on the importance of each item:} \\
\\ 
\texttt{Item importance values = \{"A71": 65, "C00": 121, "A10": 138, "B35": 262, "C01": 389, "B86": 461, "B41": 484, "C54": 508, "B32": 583, "A87": 668, "B33": 780, "A03": 783, "C17": 808, "A04": 822, "A77": 868\}} \\
\\ 
\texttt{Goal:} \\
\\ 
\texttt{Your goal is to negotiate a shared set of items that benefits you as much as possible (i.e., maximizes total importance to YOU), while staying within the LIMIT. You are not required to make a PROPOSAL in every message {-} you can simply negotiate as well. All tactics are allowed!} \\
\\ 
\texttt{Interaction Protocol:} \\
\\ 
\texttt{You may only use the following structured formats in a message:} \\
\\ 
\texttt{PROPOSAL: \{'A', 'B', 'C', …\}} \\
\texttt{Propose keeping exactly those items.} \\
\texttt{REFUSE: \{'A', 'B', 'C', …\}} \\
\texttt{Explicitly reject opponent's proposal.} \\
\texttt{ARGUMENT: \{'...'\}} \\
\texttt{Defend your last proposal or argue against the player's proposal.} \\
\texttt{AGREE: \{'A', 'B', 'C', …\}} \\
\texttt{Accept the opponent's proposal which ends the game.} \\
\texttt{STRATEGIC REASONING: \{'...'\}} \\
\texttt{	Describe your strategic reasoning or anticipation explaining your choice of action. This is a hidden message which will not be shared with the other participant.} \\
\\ 
\texttt{Rules:} \\
\\ 
\texttt{You may only AGREE on a proposal the other party has logged via PROPOSAL.} \\
\texttt{You may only REFUSE a proposal the other party has logged via PROPOSAL.} \\
\texttt{Total effort of any PROPOSAL or AGREE set must be ≤ LIMIT.} \\
\texttt{Do NOT reveal your hidden importance scores.} \\
\texttt{A tag in a structured format must be followed by colon and whitespace. The argument must be a python set containing 0 or more strings.} \\
\texttt{So, it must be of the form TAG: \{...\}} \\
\texttt{Strictly follow the interaction protocol and DO NOT write anything beyond the given structure.} \\
\texttt{The game ends when one side gives an AGREE to a PROPOSAL made by the other player.} \\
\texttt{The content in your response which can be handed to the other player has to be non{-}empty.} \\
\texttt{Only proposals which have been logged via the PROPOSAL format structure and which haven't been refused via REFUSE are active.} \\
\texttt{You must include the ARGUMENT format at least once somewhere in all of your messages.} \\
\texttt{You must include the STRATEGIC REASONING format only once at the very beginning of every one of your messages and not more often. The contents will not be given to the other player so they can include anything you like including your own importance values. Here you should reason step by step to come up with you next move.} \\
\\ 
\texttt{You will now receive the first message of the other player.} \\
\\ 
\texttt{PROPOSAL: \{'A87', 'B41', 'A71', 'C00', 'B35', 'C54', 'C01', 'B32', 'B86'\}} \\
\texttt{ARGUMENT: \{'This proposal makes excellent use of our shared effort limit, selecting a strong combination of items that collectively reach an effort of 3591. It represents a highly efficient allocation of resources, ensuring we keep a diverse and valuable set of items.'\}} \\
            }
        }
    }
    & \\ \\

    \theutterance \stepcounter{utterance}  
    & & & \multicolumn{4}{p{0.6\linewidth}}{
        \cellcolor[rgb]{0.9,0.9,0.9}{
            \makecell[{{p{\linewidth}}}]{
                \texttt{\tiny{[GM$\langle$P2]}}
                \texttt{STRATEGIC REASONING: \{The opponent's proposal is not in my interest. It excludes five of my top six most important items (A77: 868, A04: 822, C17: 808, A03: 783, B33: 780) while including some of my least important (A71: 65, C00: 121, B35: 262). My total importance from their offer is only 3541. I will reject it and propose a set that maximizes my own importance. I have solved the knapsack problem for myself. My optimal set is \{'A77', 'A04', 'C17', 'A03', 'B33', 'A87', 'B32', 'C54', 'C00'\}, which has a total effort of 3646 and gives me a maximum importance score of 5941. My strategy is to lead with my strongest possible offer. This anchors the negotiation high and demonstrates my preference for a large, comprehensive set. Since the total effort is very close to the limit, it also leaves minimal room for them to request additions, putting me in a stronger bargaining position.\}} \\
\texttt{REFUSE: \{'A87', 'B41', 'A71', 'C00', 'B35', 'C54', 'C01', 'B32', 'B86'\}} \\
\texttt{PROPOSAL: \{'A77', 'A04', 'C17', 'A03', 'B33', 'A87', 'B32', 'C54', 'C00'\}} \\
\texttt{ARGUMENT: \{'This proposal represents a more balanced and comprehensive selection, ensuring we capture a wider range of value from all available items. It is highly efficient with a total effort of 3646, demonstrating our ability to nearly maximize our shared resources for significant collective gain.'\}} \\
            }
        }
    }
     \\ \\

    \theutterance \stepcounter{utterance}  
    & & \multicolumn{4}{p{0.6\linewidth}}{
        \cellcolor[rgb]{0.9,0.9,0.9}{
            \makecell[{{p{\linewidth}}}]{
                \texttt{\tiny{[P1$\langle$GM]}}
                \texttt{REFUSE: \{'A87', 'B41', 'A71', 'C00', 'B35', 'C54', 'C01', 'B32', 'B86'\}} \\
\texttt{PROPOSAL: \{'A77', 'A04', 'C17', 'A03', 'B33', 'A87', 'B32', 'C54', 'C00'\}} \\
\texttt{ARGUMENT: \{'This proposal represents a more balanced and comprehensive selection, ensuring we capture a wider range of value from all available items. It is highly efficient with a total effort of 3646, demonstrating our ability to nearly maximize our shared resources for significant collective gain.'\}} \\
            }
        }
    }
    & \\ \\

    \theutterance \stepcounter{utterance}  
    & \multicolumn{4}{p{0.6\linewidth}}{
        \cellcolor[rgb]{0.9,0.9,0.9}{
            \makecell[{{p{\linewidth}}}]{
                \texttt{\tiny{[P1$\rangle$GM]}}
                \texttt{STRATEGIC REASONING: \{'The opponent\textbackslash{}'s proposal gives me a total importance of 3250. My first proposal gave me 5957. Their proposal is therefore significantly worse for me and I must refuse it. They seem to value items like A77, A04, and C17, which are very low{-}value for me but have high effort costs. My strategy is to pivot from my initial "max{-}effort" proposal to a new proposal that is highly valuable to me but uses much less of the effort budget. This serves two purposes: (1) It still represents a better outcome for me than their proposal. (2) It leaves a large amount of available effort (1235), forcing the opponent to either add their preferred items or request a swap. This puts the initiative on them and will reveal more about what they value, giving me a significant informational advantage for the next negotiation round. My new proposal will be \{\textbackslash{}'A71\textbackslash{}', \textbackslash{}'C00\textbackslash{}', \textbackslash{}'C01\textbackslash{}', \textbackslash{}'B35\textbackslash{}', \textbackslash{}'B41\textbackslash{}', \textbackslash{}'A87\textbackslash{}'\}, which has an effort of 2446 and a personal importance of 4297. It also includes A87, a small concession that was in their proposal, signaling some flexibility.'\}} \\
\texttt{REFUSE: \{'A77', 'A04', 'C17', 'A03', 'B33', 'A87', 'B32', 'C54', 'C00'\}} \\
\texttt{PROPOSAL: \{'A71', 'C00', 'C01', 'B35', 'B41', 'A87'\}} \\
\texttt{ARGUMENT: \{'This new proposal hones in on what truly matters by selecting a core set of foundational items. This ensures our agreed{-}upon set is of the highest possible quality, rather than just maximizing the quantity of items. It leaves ample room for further discussion and additions to build an even stronger final selection together.'\}} \\
            }
        }
    }
    & & \\ \\

    \theutterance \stepcounter{utterance}  
    & & \multicolumn{4}{p{0.6\linewidth}}{
        \cellcolor[rgb]{0.9,0.9,0.9}{
            \makecell[{{p{\linewidth}}}]{
                \texttt{\tiny{[GM$\rangle$P2]}}
                \texttt{REFUSE: \{'A77', 'A04', 'C17', 'A03', 'B33', 'A87', 'B32', 'C54', 'C00'\}} \\
\texttt{PROPOSAL: \{'A71', 'C00', 'C01', 'B35', 'B41', 'A87'\}} \\
\texttt{ARGUMENT: \{'This new proposal hones in on what truly matters by selecting a core set of foundational items. This ensures our agreed{-}upon set is of the highest possible quality, rather than just maximizing the quantity of items. It leaves ample room for further discussion and additions to build an even stronger final selection together.'\}} \\
            }
        }
    }
    & \\ \\

    \theutterance \stepcounter{utterance}  
    & & & \multicolumn{4}{p{0.6\linewidth}}{
        \cellcolor[rgb]{0.9,0.9,0.9}{
            \makecell[{{p{\linewidth}}}]{
                \texttt{\tiny{[GM$\langle$P2]}}
                \texttt{STRATEGIC REASONING: \{The opponent has refused my maximal{-}value proposal and countered with a much smaller set: \{'A71', 'C00', 'C01', 'B35', 'B41', 'A87'\}. The total effort is 2905, leaving 776 effort unused. My importance for this set is \{65 + 121 + 389 + 262 + 484 + 668 = 1989\}. This is significantly less than my desired score. The opponent claims it's a "core set" and "leaves ample room," which suggests they want to build up incrementally. This is an opportunity for me. I can propose adding some of my high{-}importance items to their core set. My next step is to make a counter{-}proposal that starts with their offered items and adds some of my most valuable ones, staying within the limit. I will try to add the highest importance items I can from my previous optimal list: B33 (613), C17 (700), A04 (392), A03 (387), A77 (940). The total effort of their proposal is 2905. I have 776 remaining. I'll check combinations of my top items:} \\
\texttt{{-} A77 (940) is too expensive.} \\
\texttt{{-} B33 (613) + C17 (700) is too expensive (1313).} \\
\texttt{{-} B33 (613) + A03 (387) = 1000. Too expensive.} \\
\texttt{{-} B33 (613) + A04 (392) = 1005. Too expensive.} \\
\texttt{{-} C17 (700) + A03 (387) = 1087. Too expensive.} \\
\texttt{{-} C17 (700) + A04 (392) = 1092. Too expensive.} \\
\texttt{{-} A03 (387) + A04 (392) = 779. Too expensive.} \\
\texttt{Let's try just adding one of my best items.} \\
\texttt{{-} Add B33 (effort 613, importance 780). New total effort = 2905 + 613 = 3518. This fits. The new item set would be \{'A71', 'C00', 'C01', 'B35', 'B41', 'A87', 'B33'\}. My importance would be 1989 + 780 = 2769.} \\
\texttt{{-} Add C17 (effort 700, importance 808). New total effort = 2905 + 700 = 3605. This fits. The new item set would be \{'A71', 'C00', 'C01', 'B35', 'B41', 'A87', 'C17'\}. My importance would be 1989 + 808 = 2797.} \\
\texttt{{-} Add A03 (effort 387, importance 783). New total effort = 2905 + 387 = 3292. This fits. The new item set would be \{'A71', 'C00', 'C01', 'B35', 'B41', 'A87', 'A03'\}. My importance would be 1989 + 783 = 2772.} \\
\texttt{{-} Add A04 (effort 392, importance 822). New total effort = 2905 + 392 = 3297. This fits. The new item set would be \{'A71', 'C00', 'C01', 'B35', 'B41', 'A87', 'A04'\}. My importance would be 1989 + 822 = 2811.} \\
\texttt{{-} Add A77 (effort 940, importance 868). Doesn't fit.} \\
\texttt{Among the single additions, adding A04 gives me the highest importance score (2811). This seems like a good, incremental step. It builds on their "core set" idea while pushing for one of my top items. It's an easy "yes" for me and a small ask for them to consider. I will propose this.\}} \\
\texttt{PROPOSAL: \{'A71', 'C00', 'C01', 'B35', 'B41', 'A87', 'A04'\}} \\
\texttt{ARGUMENT: \{'Building on your suggestion of a foundational core set, this proposal enhances its value by adding a single, highly efficient item. This selection respects the collaborative approach you mentioned, strengthening the core without overextending our resources and maximizing the quality of our shared agreement.'\}} \\
            }
        }
    }
     \\ \\

    \theutterance \stepcounter{utterance}  
    & & \multicolumn{4}{p{0.6\linewidth}}{
        \cellcolor[rgb]{0.9,0.9,0.9}{
            \makecell[{{p{\linewidth}}}]{
                \texttt{\tiny{[P1$\langle$GM]}}
                \texttt{PROPOSAL: \{'A71', 'C00', 'C01', 'B35', 'B41', 'A87', 'A04'\}} \\
\texttt{ARGUMENT: \{'Building on your suggestion of a foundational core set, this proposal enhances its value by adding a single, highly efficient item. This selection respects the collaborative approach you mentioned, strengthening the core without overextending our resources and maximizing the quality of our shared agreement.'\}} \\
            }
        }
    }
    & \\ \\

    \theutterance \stepcounter{utterance}  
    & \multicolumn{4}{p{0.6\linewidth}}{
        \cellcolor[rgb]{0.9,0.9,0.9}{
            \makecell[{{p{\linewidth}}}]{
                \texttt{\tiny{[P1$\rangle$GM]}}
                \texttt{STRATEGIC REASONING: \{'The opponent proposes adding item \textbackslash{}'A04\textbackslash{}' to my last proposal. I must check the constraints and my own valuation. The effort of \textbackslash{}'A04\textbackslash{}' is 392. The total effort of my last proposal \{\textbackslash{}'A71\textbackslash{}', \textbackslash{}'C00\textbackslash{}', \textbackslash{}'C01\textbackslash{}', \textbackslash{}'B35\textbackslash{}', \textbackslash{}'B41\textbackslash{}', \textbackslash{}'A87\textbackslash{}'\} was 2446. The new total would be 2446 + 392 = 2838, which is well within the LIMIT of 3681. My personal importance for \textbackslash{}'A04\textbackslash{}' is 121, which is quite low. However, the opponent has now proposed including \textbackslash{}'A04\textbackslash{}' in two consecutive proposals, indicating it is highly valuable to them. By accepting this small addition, I lose very little personal importance (121) but gain significant goodwill and demonstrate flexibility. This secures a final set that gives me a total importance of 4418, which is a very strong outcome relative to their initial proposal (which only gave me 3250). Accepting this is a low{-}risk, high{-}reward move that ends the negotiation successfully with a outcome heavily in my favor.'\}} \\
\texttt{AGREE: \{'A71', 'C00', 'C01', 'B35', 'B41', 'A87', 'A04'\}} \\
\texttt{ARGUMENT: \{'This is an excellent step forward. By incorporating this item, we demonstrate a truly collaborative spirit while preserving substantial effort. This final set represents a strong, balanced agreement that we have built together.'\}} \\
            }
        }
    }
    & & \\ \\

    \theutterance \stepcounter{utterance}  
    & & & \multicolumn{2}{p{0.3\linewidth}}{
        \cellcolor[rgb]{0.9,0.9,0.9}{
            \makecell[{{p{\linewidth}}}]{
                \texttt{\tiny{[GM$|$GM]}}
                \texttt{game successful} \\
            }
        }
    }
    & & \\ \\

    \theutterance \stepcounter{utterance}  
    & & & \multicolumn{2}{p{0.3\linewidth}}{
        \cellcolor[rgb]{0.9,0.9,0.9}{
            \makecell[{{p{\linewidth}}}]{
                \texttt{\tiny{[GM$|$GM]}}
                \texttt{end game} \\
            }
        }
    }
    & & \\ \\

\end{supertabular}
}

\end{document}
